
% =====================================================================
\begin{thebibliography}{999}\small
\setlength{\itemsep}{1pt}
\bibitem{feynman1982} R.~P.~Feynman, ``Simulating physics with computers,'' Int.\ J.\ Theor.\ Phys.\ 21, 467--488 (1982). \url{https://doi.org/10.1007/BF02650179}
\bibitem{lloyd1996} S.~Lloyd, ``Universal quantum simulators,'' Science 273, 1073--1078 (1996). \url{https://doi.org/10.1126/science.273.5278.1073}
\bibitem{aspuru2005} A.~Aspuru-Guzik, A.~D.~Dutoi, P.~J.~Love, M.~Head-Gordon, ``Simulated quantum computation of molecular energies,'' Science 309, 1704--1707 (2005). \url{https://doi.org/10.1126/science.1113479}
\bibitem{reiher2017} M.~Reiher, N.~Wiebe, K.~M.~Svore, D.~Wecker, M.~Troyer, ``Elucidating reaction mechanisms on quantum computers,'' PNAS 114, 7555--7560 (2017). \url{https://doi.org/10.1073/pnas.1619152114}
\bibitem{bauer2020} B.~Bauer, S.~Bravyi, M.~Motta, G.~K.-L.~Chan, ``Quantum algorithms for quantum chemistry and quantum materials science,'' Chem.\ Rev.\ 120, 12685--12717 (2020). \url{https://doi.org/10.1021/acs.chemrev.9b00829}
\bibitem{white1992} S.~R.~White, ``Density matrix formulation for quantum renormalization groups,'' Phys.\ Rev.\ Lett.\ 69, 2863 (1992). \url{https://doi.org/10.1103/PhysRevLett.69.2863}
\bibitem{chan2011} G.~K.-L.~Chan, S.~Sharma, ``The density matrix renormalization group in quantum chemistry,'' Annu.\ Rev.\ Phys.\ Chem.\ 62, 465--481 (2011). \url{https://doi.org/10.1146/annurev-physchem-032210-103338}
\bibitem{carleo2017} G.~Carleo, M.~Troyer, ``Solving the quantum many-body problem with artificial neural networks,'' Science 355, 602--606 (2017). \url{https://doi.org/10.1126/science.aag2302}
\bibitem{carleo2019} G.~Carleo et al., ``Machine learning and the physical sciences,'' Rev.\ Mod.\ Phys.\ 91, 045002 (2019). \url{https://doi.org/10.1103/RevModPhys.91.045002}
\bibitem{jumper2021} J.~Jumper et al., ``Highly accurate protein structure prediction with AlphaFold,'' Nature 596, 583--589 (2021). \url{https://doi.org/10.1038/s41586-021-03819-2}
\bibitem{arute2019} F.~Arute et al., ``Quantum supremacy using a programmable superconducting processor,'' Nature 574, 505--510 (2019). \url{https://doi.org/10.1038/s41586-019-1666-5}
\bibitem{zhong2020} H.-S.~Zhong et al., ``Quantum computational advantage using photons,'' Science 370, 1460--1463 (2020). \url{https://doi.org/10.1126/science.abe8770}
\bibitem{kim2023} Y.~Kim et al., ``Evidence for the utility of quantum computing before fault tolerance,'' Nature 618, 500--505 (2023). \url{https://doi.org/10.1038/s41586-023-06096-3}
\bibitem{semeghini2021} G.~Semeghini et al., ``Probing topological spin liquids on a programmable quantum simulator,'' Science 374, 1242--1247 (2021). \url{https://doi.org/10.1126/science.abi8794}
\bibitem{preskill2018} J.~Preskill, ``Quantum computing in the NISQ era and beyond,'' Quantum 2, 79 (2018). \url{https://doi.org/10.22331/q-2018-08-06-79}
\bibitem{bharti2022} K.~Bharti et al., ``Noisy intermediate-scale quantum algorithms,'' Rev.\ Mod.\ Phys.\ 94, 015004 (2022). \url{https://doi.org/10.1103/RevModPhys.94.015004}
\bibitem{aaronson2015} S.~Aaronson, ``Read the fine print,'' Nat.\ Phys.\ 11, 291--293 (2015). \url{https://doi.org/10.1038/nphys3272}
\bibitem{daley2022} A.~J.~Daley et al., ``Practical quantum advantage in quantum simulation,'' Nature 607, 667--676 (2022). \url{https://doi.org/10.1038/s41586-022-04940-6}
\bibitem{dalzell2023} A.~M.~Dalzell et al., ``Quantum algorithms: A survey of applications and end-to-end complexities,'' arXiv:2310.03011 (2023). \url{https://arxiv.org/abs/2310.03011}
\bibitem{gent2024} E.~Gent, ``Why AI could eat quantum computing's lunch,'' MIT Technology Review, 7 Nov 2024. \url{https://www.technologyreview.com/2024/11/07/1106730/why-ai-could-eat-quantum-computings-lunch/}
\bibitem{babbush2025} R.~Babbush et al., ``The grand challenge of quantum applications,'' arXiv:2511.09124 (2025). \url{https://arxiv.org/abs/2511.09124}
\bibitem{hoefler2023} T.~Hoefler, T.~H\"aner, M.~Troyer, ``Disentangling hype from practicality: on realistically achieving quantum advantage,'' Commun.\ ACM 66(5):82--87 (2023). \url{https://arxiv.org/abs/2307.00523}
\bibitem{babbush2021} R.~Babbush, J.~R.~McClean, M.~Newman, C.~Gidney, S.~Boixo, H.~Neven, ``Focus beyond quadratic speedups for error-corrected quantum advantage,'' PRX Quantum 2, 010103 (2021). \url{https://arxiv.org/abs/2011.04149}
\bibitem{beverland2022} M.~E.~Beverland et al., ``Assessing requirements to scale to practical quantum advantage,'' arXiv:2211.07629v1 (2022), Tables~I and~IV and Appendix~F. \url{https://arxiv.org/abs/2211.07629}
\bibitem{beverland_qualtran} Qualtran documentation, ``Beverland et al.\ model'' (reproduction of Appendix F of~\cite{beverland2022}). \url{https://qualtran.readthedocs.io/en/latest/surface_code/beverland_et_al_model.html}
\bibitem{hangleiter2026} D.~Hangleiter, ``Has quantum advantage been achieved?'' arXiv:2603.09901 (2026). \url{https://arxiv.org/abs/2603.09901}
\bibitem{pan2022} F.~Pan, K.~Chen, P.~Zhang, ``Solving the sampling problem of the Sycamore quantum circuits,'' Phys.\ Rev.\ Lett.\ 129, 090502 (2022). \url{https://arxiv.org/abs/2111.03011}
\bibitem{zhong2021} H.-S.~Zhong et al., ``Phase-programmable Gaussian boson sampling using stimulated squeezed light,'' Phys.\ Rev.\ Lett.\ 127, 180502 (2021). \url{https://arxiv.org/abs/2106.15534}

\bibitem{villalonga2021} B.~Villalonga, M.~Y.~Niu, L.~Li, H.~Neven, J.~C.~Platt, V.~N.~Smelyanskiy, S.~Boixo, ``Efficient approximation of experimental Gaussian boson sampling,'' arXiv:2109.11525 (2021). \url{https://arxiv.org/abs/2109.11525}

\bibitem{martinez2023} J.~Mart\'inez-Cifuentes, K.~M.~Fonseca-Romero, N.~Quesada, ``Classical models may be a better explanation of the Jiuzhang 1.0 Gaussian boson sampler than its targeted squeezed light model,'' Quantum 7, 1076 (2023). \url{https://arxiv.org/abs/2207.10058}

\bibitem{madsen2022} L.~S.~Madsen et al., ``Quantum computational advantage with a programmable photonic processor,'' Nature 606, 75--81 (2022). \url{https://doi.org/10.1038/s41586-022-04725-x}

\bibitem{deng2023} Y.-H.~Deng et al., ``Gaussian boson sampling with pseudo-photon-number-resolving detectors and quantum computational advantage,'' Phys.\ Rev.\ Lett.\ 131, 150601 (2023). \url{https://arxiv.org/abs/2304.12240}

\bibitem{anand2026} K.~Anand, Z.~Chen, M.~Cryan, G.~Freifeld, L.~A.~Goldberg, H.~Guo, X.~Zhang, ``Simulating Gaussian boson sampling on graphs in polynomial time,'' Quantum 10, 2268 (2026). \url{https://doi.org/10.22331/q-2026-11-10-2268}

\bibitem{qsub2025} ``qsub: quantum subroutine resource-estimation package,'' repository of the USC Information Sciences Institute, commit \texttt{cc8d3ce}, files under \texttt{QRE/CFD/2025-06}. \url{https://github.com/isi-usc-edu/qsub}

\bibitem{oh2024} C.~Oh, M.~Liu, Y.~Alexeev, B.~Fefferman, L.~Jiang, ``Classical algorithm for simulating experimental Gaussian boson sampling,'' Nat.\ Phys.\ 20, 1461 (2024). \url{https://arxiv.org/abs/2306.03709}
\bibitem{tindall2024} J.~Tindall, M.~Fishman, E.~M.~Stoudenmire, D.~Sels, ``Efficient tensor network simulation of IBM's Eagle kicked Ising experiment,'' PRX Quantum 5, 010308 (2024). \url{https://arxiv.org/abs/2306.14887}
\bibitem{aps2024} APS Physics 17, 13 (2024). \url{https://physics.aps.org/articles/v17/13}
\bibitem{begusic2023} T.~Begu\v{s}i\'c, G.~K.-L.~Chan, ``Fast classical simulation of evidence for the utility of quantum computing before fault tolerance,'' arXiv:2306.16372 (2023). \url{https://arxiv.org/abs/2306.16372}
\bibitem{begusic2024} T.~Begu\v{s}i\'c, J.~Gray, G.~K.-L.~Chan, ``Fast and converged classical simulations of evidence for the utility of quantum computing before fault tolerance,'' Sci.\ Adv.\ 10, eadk4321 (2024). \url{https://arxiv.org/abs/2308.05077}
\bibitem{dallatorre} E.~G.~Dalla Torre, M.~M.~Roses, ``Dissipative mean-field theory of IBM utility experiment,'' arXiv:2308.01339. \url{https://arxiv.org/abs/2308.01339}
\bibitem{king2025} A.~D.~King et al., ``Beyond-classical computation in quantum simulation,'' Science 388, 199 (2025). \url{https://arxiv.org/abs/2403.00910}
\bibitem{tindall2025} J.~Tindall et al., ``Dynamics of disordered quantum systems with two- and three-dimensional tensor networks,'' Science 392, 868--872 (2026); arXiv:2503.05693 (first posted 7 March 2025). \url{https://arxiv.org/abs/2503.05693}
\bibitem{mauron2025} L.~Mauron, G.~Carleo, ``Challenging the quantum advantage frontier with large-scale classical simulations of annealing dynamics,'' arXiv:2503.08247 (2025). Jastrow--Feenberg ansatz. \url{https://arxiv.org/abs/2503.08247}
\bibitem{kingreply} A.~D.~King et al., ``Comment on arXiv:2503.05693,'' arXiv:2504.06283 (2025). \url{https://arxiv.org/abs/2504.06283}
\bibitem{mauron2025tee} L.~Mauron, Z.~Denis, J.~Nys, G.~Carleo, ``Predicting topological entanglement entropy in a Rydberg analog simulator,'' Nat.\ Phys.\ 21, 1332 (2025). \url{https://arxiv.org/abs/2406.19872}
\bibitem{brydges2019} T.~Brydges et al., ``Probing R\'enyi entanglement entropy via randomized measurements,'' Science 364, 260 (2019). \url{https://arxiv.org/abs/1806.05747}
\bibitem{haghshenas2025} R.~Haghshenas et al., ``Digital quantum magnetism on a trapped-ion quantum computer,'' arXiv:2503.20870v3 (2026). \url{https://arxiv.org/abs/2503.20870}
\bibitem{google2025otoc} Google Quantum AI and Collaborators, ``Observation of constructive interference at the edge of quantum ergodicity,'' Nature 646, 825--830 (2025), doi:10.1038/s41586-025-09526-6; preprint arXiv:2506.10191. \url{https://arxiv.org/abs/2506.10191}
\bibitem{googleblog2025} Google Research, ``A verifiable quantum advantage,'' 22 Oct 2025 (source of the 2.1\,h and 3.2\,y figures). \url{https://research.google/blog/a-verifiable-quantum-advantage/}
\bibitem{bermejo2026} P.~Bermejo, B.~Villalonga, B.~Ware, G.~Vidal, A.~Szasz, ``Tensor networks with belief propagation cannot feasibly simulate Google's quantum echoes experiment,'' arXiv:2604.15427 (2026). \url{https://arxiv.org/abs/2604.15427}
\bibitem{barron2026} S.~V.~Barron, B.~Mitchell, V.~Tripathi et al., ``Observable estimation in the absence of classical verification,'' arXiv:2607.25998 (2026). \url{https://arxiv.org/abs/2607.25998}
\bibitem{ibmalgo2026} IBM Newsroom, announcement of IBM--Algorithmiq operator Loschmidt echo results, 30 Jul 2026 (context only). \url{https://newsroom.ibm.com/2026-07-30-ibm-and-algorithmiq-demonstrate-quantum-advantage,-establishing-a-framework-for-trusted-quantum-computation-beyond-classical-verification}
\bibitem{pq_ibm3} PostQuantum, ``IBM declares a quantum advantage era'' (2026), on the Quantum Advantage Tracker. \url{https://postquantum.com/industry-news/ibm-three-quantum-advantage-papers/}
\bibitem{granet2025} E.~Granet, S.-H.~Lin, K.~H\'emery, R.~Haghshenas et al., ``Superconducting pairing correlations on a trapped-ion quantum computer,'' arXiv:2511.02125v3 (2026). \url{https://arxiv.org/abs/2511.02125}
\bibitem{yoshioka2024} N.~Yoshioka, T.~Okubo, Y.~Suzuki, Y.~Koizumi, W.~Mizukami, ``Hunting for quantum-classical crossover in condensed matter problems,'' npj Quantum Inf.\ 10 (2024); arXiv:2210.14109v3, Supplementary Tables S1, S2 and S12. \url{https://arxiv.org/abs/2210.14109}
\bibitem{lee2023} S.~Lee et al., ``Evaluating the evidence for exponential quantum advantage in ground-state quantum chemistry,'' Nat.\ Commun.\ 14, 1952 (2023). \url{https://www.nature.com/articles/s41467-023-37587-6}
\bibitem{zhai2026} H.~Zhai, C.~Li, X.~Zhang, Z.~Li, S.~Lee, G.~K.-L.~Chan, ``Classical computational simulation of the FeMo-cofactor model to chemical accuracy and its implications,'' arXiv:2601.04621v2 (2026). \url{https://arxiv.org/abs/2601.04621}
\bibitem{lee2021thc} J.~Lee et al., ``Even more efficient quantum computations of chemistry through tensor hypercontraction,'' PRX Quantum 2, 030305 (2021); arXiv:2011.03494v3, Tables~IV and~V and Sec.~IV. \url{https://arxiv.org/abs/2011.03494}
\bibitem{reinholdt2025} P.~Reinholdt et al., ``Critical limitations in quantum-selected configuration interaction methods,'' J.\ Chem.\ Theory Comput.\ 21, 6811--6822 (2025); arXiv:2501.07231v5, Secs.~4.1 and 4.3. \url{https://doi.org/10.1021/acs.jctc.5c00375}
\bibitem{shirakawa2025} T.~Shirakawa et al., ``Closed-loop calculations of electronic structure on a quantum processor and a classical supercomputer at full scale,'' arXiv:2511.00224 (first submitted 31 October 2025). \url{https://arxiv.org/abs/2511.00224}
\bibitem{bonilla2026} N.~Bonilla Vargas, ``Machine learning for sample-based quantum diagonalization: a review,'' arXiv:2608.05314 (2026). \url{https://arxiv.org/abs/2608.05314}
\bibitem{luise2025} G.~Luise et al., ``Accurate and scalable exchange-correlation with deep learning,'' arXiv:2506.14665 (2025). \url{https://arxiv.org/abs/2506.14665}
\bibitem{skala10} Microsoft Research, ``Skala 1.0 model card.'' \url{https://microsoft.github.io/skala/model-card/skala-1.0.html}
\bibitem{skala11} Microsoft Research, ``Skala 1.1 model card.'' \url{https://microsoft.github.io/skala/model-card/skala-1.1.html}
\bibitem{goings2022} J.~J.~Goings et al., ``Reliably assessing the electronic structure of cytochrome P450 on today's classical computers and tomorrow's quantum computers,'' PNAS 119, e2203533119 (2022). \url{https://pmc.ncbi.nlm.nih.gov/articles/PMC9499570}
\bibitem{low2025sa} G.~H.~Low, R.~King, D.~W.~Berry, Q.~Han, A.~E.~DePrince III, A.~F.~White, R.~Babbush, R.~D.~Somma, N.~C.~Rubin, ``Fast quantum simulation of electronic structure by spectral amplification,'' Phys.\ Rev.\ X 15, 041016 (2025), doi:10.1103/pb2g-j9cw. Earlier version read for the gate counts: arXiv:2502.15882v1, titled ``\ldots by spectrum amplification''. \url{https://arxiv.org/abs/2502.15882}
\bibitem{low2026dense} G.~H.~Low, W.~J.~Huggins, D.~W.~Berry, T.~Khattar, A.~F.~White, N.~C.~Rubin, R.~Babbush, ``A denser planar surface code,'' arXiv:2605.30455v1 (2026), Fig.~1 and Appendix~F, Table~XXVII. \url{https://arxiv.org/abs/2605.30455}
\bibitem{fomichev2025} S.~Fomichev, P.~A.~M.~Casares, J.~Soni et al., ``Fast simulations of X-ray absorption spectroscopy for battery materials on a quantum computer,'' arXiv:2506.15784v3 (5 November 2026), Table~1 and Introduction; version 2 (20 August 2025) for the classical estimate. \url{https://arxiv.org/abs/2506.15784}
\bibitem{ibmroadmap2025} IBM Newsroom, 10 Jun 2025. \url{https://newsroom.ibm.com/2025-06-10-IBM-Sets-the-Course-to-Build-Worlds-First-Large-Scale,-Fault-Tolerant-Quantum-Computer-at-New-IBM-Quantum-Data-Center}
\bibitem{hpcwire2025} HPCwire, ``IBM sets 2029 target for fault-tolerant quantum computing,'' 10 Jun 2025. \url{https://www.hpcwire.com/2025/06/10/ibm-sets-2029-target-for-fault-tolerant-quantum-computing/}
\bibitem{quantinuum2024} H.~Buhrman, C.~Langer, Quantinuum technical perspective, 17 Sep 2024. \url{https://www.quantinuum.com/blog/technical-perspective-by-the-end-of-the-decade-we-will-deliver-universal-fault-tolerant-quantum-computing}
\bibitem{quantFiling} Quantinuum Inc., Form DRS (2026), technology roadmap section. \url{https://www.sec.gov/Archives/edgar/data/0002110105/000162827926000199/filename1.htm}
\bibitem{merchant2023} A.~Merchant, S.~Batzner, S.~S.~Schoenholz, M.~Aykol, G.~Cheon, E.~D.~Cubuk, ``Scaling deep learning for materials discovery,'' Nature 624, 80 (2023). \url{https://doi.org/10.1038/s41586-023-06735-9}
\bibitem{robert2021} A.~Robert, P.~K.~Barkoutsos, S.~Woerner, I.~Tavernelli, ``Resource-efficient quantum algorithm for protein folding,'' npj Quantum Inf.\ 7, 38 (2021). \url{https://arxiv.org/abs/1908.02163}
\bibitem{doga2024} H.~Doga et al., ``A perspective on protein structure prediction using quantum computers,'' J.\ Chem.\ Theory Comput.\ 20, 3359--3378 (2024); arXiv:2312.00875 (discussion of Fig.~9). \url{https://doi.org/10.1021/acs.jctc.4c00067}
\bibitem{zhang2026} Y.~Zhang, Y.~Yang, W.~Martin et al., ``A quantum framework for protein binding-site structure prediction on utility-level quantum processors,'' Adv.\ Sci.\ (2026), doi:10.1002/advs.202513641; arXiv:2506.22677v2. \url{https://arxiv.org/abs/2506.22677}
\bibitem{kras2025} M.~Ghazi Vakili et al., ``Quantum-computing-enhanced algorithm unveils potential KRAS inhibitors,'' Nat.\ Biotechnol.\ 43, 1954 (2025). \url{https://pmc.ncbi.nlm.nih.gov/articles/PMC12700792}
\bibitem{uoft2025} University of Toronto Donnelly Centre, 22 Jan 2025. \url{https://thedonnellycentre.utoronto.ca/news/u-t-researchers-develop-new-approach-using-quantum-computers-accelerate-drug-discovery}
\bibitem{hibat2021} M.~Hibat-Allah, E.~M.~Inack, R.~Wiersema, R.~G.~Melko, J.~Carrasquilla, ``Variational neural annealing,'' Nat.\ Mach.\ Intell.\ 3, 952 (2021). \url{https://arxiv.org/abs/2101.10154}
\bibitem{hibat2021b} M.~Hibat-Allah, APS March Meeting 2021 invited talk. \url{https://meetings-archive.aps.org/mar/2021/a21/5}
\bibitem{munoz2025} H.~Munoz-Bauza, D.~A.~Lidar, ``Scaling advantage in approximate optimization with quantum annealing,'' Phys.\ Rev.\ Lett.\ 134, 160601 (2025); arXiv:2401.07184, first posted 14 January 2024. \url{https://arxiv.org/abs/2401.07184}
\bibitem{dwave_qac} D-Wave publication summary (2025). \url{https://dwavesys.com/resources/publication/scaling-advantage-in-approximate-optimization-with-quantum-annealing}
\bibitem{pawlowski2026} J.~Paw{\l}owski, P.~Tarasiuk, J.~Tuziemski, {\L}.~Pawela, B.~Gardas, ``Toward quantum scaling advantage in approximate optimization,'' Phys.\ Rev.\ Applied 26, 014024 (2026). \url{https://arxiv.org/abs/2505.22514}
\bibitem{angelini2022} M.~C.~Angelini, F.~Ricci-Tersenghi, ``Modern graph neural networks do worse than classical greedy algorithms,'' Nat.\ Mach.\ Intell.\ (2023). \url{https://arxiv.org/abs/2206.13211}
\bibitem{schuetz2023} M.~J.~A.~Schuetz, J.~K.~Brubaker, H.~G.~Katzgraber, ``Reply,'' Nat.\ Mach.\ Intell.\ 5, 32 (2023). \url{https://arxiv.org/abs/2302.03602}
\bibitem{tang2023} E.~Tang, ``Quantum machine learning without any quantum,'' PhD thesis, University of Washington (2023). \url{https://digital.lib.washington.edu/researchworks/items/16909f93-67a2-4620-aa59-0f492fd7290f/full}
\bibitem{tang2021} E.~Tang, ``Quantum principal component analysis only achieves an exponential speedup because of its state preparation assumptions,'' Phys.\ Rev.\ Lett.\ 127, 060503 (2021). \url{https://arxiv.org/abs/1811.00414}
\bibitem{liu2021} Y.~Liu, S.~Arunachalam, K.~Temme, ``A rigorous and robust quantum speed-up in supervised machine learning,'' Nat.\ Phys.\ 17, 1013 (2021). \url{https://arxiv.org/abs/2010.02174}
\bibitem{huang2022} H.-Y.~Huang, R.~Kueng, G.~Torlai, V.~V.~Albert, J.~Preskill, ``Provably efficient machine learning for quantum many-body problems,'' Science 377, eabk3333 (2022). \url{https://arxiv.org/abs/2106.12627}
\bibitem{huang2022exp} H.-Y.~Huang et al., ``Quantum advantage in learning from experiments,'' Science 376, 1182 (2022). \url{https://arxiv.org/abs/2112.00778}
\bibitem{lovato2022} A.~Lovato, C.~Adams, G.~Carleo, N.~Rocco, ``Hidden-nucleons neural-network quantum states for the nuclear many-body problem,'' Phys.\ Rev.\ Research (2022). \url{https://arxiv.org/abs/2206.10021}
\bibitem{gnech2024} A.~Gnech et al., ``Distilling the essential elements of nuclear binding via neural-network quantum states,'' Phys.\ Rev.\ Lett.\ (2024). \url{https://arxiv.org/abs/2308.16266}
\bibitem{lovato2026} A.~Lovato et al., ``Neural-network quantum states for the nuclear many-body problem,'' arXiv:2602.13826 (2026). \url{https://arxiv.org/abs/2602.13826}
\bibitem{rhodes2024} M.~L.~Rhodes, M.~Kreshchuk, S.~Pathak, ``Exponential improvements in the simulation of lattice gauge theories using near-optimal techniques,'' PRX Quantum 5, 040347 (2024). \url{https://arxiv.org/abs/2405.10416}
\bibitem{carena2022} M.~Carena, E.~Gustafson, H.~Lamm et al., ``Gauge theory couplings on anisotropic lattices,'' arXiv:2208.10417 (2022); cites the $\sim10^{49}$ T shear-viscosity estimate from earlier work. \url{https://arxiv.org/abs/2208.10417}
\bibitem{lamm2026} H.~Lamm, ``Ether of orbifolds,'' arXiv:2603.29091v3 (2026); Table of T-gate estimates for a fiducial SU(3) shear-viscosity calculation. \url{https://arxiv.org/abs/2603.29091}
\bibitem{osorio2025} S.~Osorio Perez, E.~M.~Murairi, E.~J.~Gustafson, H.~Lamm, ``Primitive quantum gates for an SU(3) discrete subgroup: $\Sigma(72\times3)$,'' arXiv:2511.17437 (2025). \url{https://arxiv.org/abs/2511.17437}
\bibitem{draper2026} P.~Draper, ``Block encoding non-Abelian lattice gauge theory,'' arXiv:2608.17115 (2026). \url{https://arxiv.org/abs/2608.17115}
\bibitem{jennings2025} D.~Jennings et al., ``An end-to-end quantum algorithm for nonlinear fluid dynamics with bounded quantum advantage,'' PRX Quantum 7, 033060 (2026); arXiv:2512.03758v2. \url{https://arxiv.org/abs/2512.03758}
\bibitem{zhuang2025} X.-N.~Zhuang et al., ``A pathway to practical quantum advantage in solving Navier--Stokes equations,'' arXiv:2509.08807 (2025). \url{https://arxiv.org/abs/2509.08807}
\bibitem{mello2025} A.~F.~Mello, M.~Collura, E.~M.~Stoudenmire, R.~Levy, ``Magic of the well,'' arXiv:2512.03177 (2025). \url{https://arxiv.org/abs/2512.03177}
\bibitem{chen2024} C.~Chen, D.~T.~Nguyen, S.~J.~Lee, N.~A.~Baker, A.~S.~Karakoti, L.~Lauw, C.~Owen et al., ``Accelerating computational materials discovery with machine learning and cloud high-performance computing: from large-scale screening to experimental validation,'' J.\ Am.\ Chem.\ Soc.\ 146(29), 20009 (2024), doi:10.1021/jacs.4c03849; arXiv:2401.04070. \url{https://doi.org/10.1021/jacs.4c03849}
\bibitem{belagali2026} H.~Belagali et al., ``Efficient classical simulation of large-scale unitary cluster Jastrow circuits,'' arXiv:2607.21337 (2026). \url{https://arxiv.org/abs/2607.21337}
\bibitem{ciavarella2025} A.~N.~Ciavarella, I.~M.~Burbano, C.~W.~Bauer, ``Efficient truncations of SU($N_c$) lattice gauge theory for quantum simulation,'' arXiv:2503.11888v5 (2026). \url{https://arxiv.org/abs/2503.11888}
\bibitem{hirsch1985} J.~E.~Hirsch, ``Two-dimensional Hubbard model: numerical simulation study,'' Phys.\ Rev.\ B 31, 4403 (1985). \url{https://doi.org/10.1103/PhysRevB.31.4403}
\bibitem{leblanc2015} J.~P.~F.~LeBlanc et al.\ (Simons Collaboration on the Many-Electron Problem), ``Solutions of the two-dimensional Hubbard model: benchmarks and results from a wide range of numerical algorithms,'' Phys.\ Rev.\ X 5, 041041 (2015). \url{https://doi.org/10.1103/PhysRevX.5.041041}
\bibitem{berry2024} D.~W.~Berry et al., ``Rapid initial state preparation for the quantum simulation of strongly correlated molecules,'' arXiv:2409.11748v1 (2024); abstract: FeMoco estimate of $7.3\times10^{10}$ Toffoli including preparation. \url{https://arxiv.org/abs/2409.11748}
\bibitem{leviatan2026} E.~Leviatan et al., ``Resolving structure in prethermal Floquet dynamics with precision quantum computation,'' arXiv:2607.24937v1 (27 July 2026). \url{https://arxiv.org/abs/2607.24937}
\bibitem{martiel2026} S.~Martiel, J.-U.~Chung, A.~Seif, S.~Ghosh, I.~Hincks, A.~Deshpande, B.~Fefferman, J.~M.~Gambetta, A.~Javadi-Abhari, ``Sampling hard circuits with verifiably high fidelity,'' arXiv:2607.25941 (v1 28 July 2026; v3 2 September 2026). \url{https://arxiv.org/abs/2607.25941}
\bibitem{manabe2026} H.~Manabe, H.~Gu, F.~Pan, ``Classical simulation and design frontiers for IBM's doped Clifford sampling experiment,'' arXiv:2608.13110 (v1 13 August 2026; v2 29 October 2026). \url{https://arxiv.org/abs/2608.13110}
\bibitem{zheng2026} Y.-G.~Zheng, Y.-C.~Shen, W.-Y.~Zhang et al., ``Quantifying quantum computational advantage on a processor of ultracold atoms,'' Phys.\ Rev.\ X 16, 021057 (2026); arXiv:2210.08556v2 (11 April 2026; version 1 of 16 October 2022 has another title and a 32-site system). \url{https://doi.org/10.1103/5mqb-3by1}
\bibitem{ebadi2022} S.~Ebadi et al., ``Quantum optimization of maximum independent set using Rydberg atom arrays,'' Science 376, 1209--1215 (2022); arXiv:2202.09372. \url{https://doi.org/10.1126/science.abo6587}
\bibitem{andrist2023} R.~S.~Andrist et al., ``Hardness of the maximum-independent-set problem on unit-disk graphs and prospects for quantum speedups,'' Phys.\ Rev.\ Research 5, 043277 (2023); arXiv:2307.09442. \url{https://doi.org/10.1103/PhysRevResearch.5.043277}
\bibitem{merz2026} K.~M.~Merz,~Jr., et al., ``Crossing the 12,000-atom barrier with heterogeneous quantum-classical supercomputing: quantum chemistry of protein-ligand complexes,'' arXiv:2605.01138v2 (5 September 2026). \url{https://arxiv.org/abs/2605.01138}
\bibitem{legeza2026} \"O.~Legeza, A.~Menczer, M.~A.~Werner, S.~S.~Xantheas, F.~Neese, M.~Ganahl, C.~Brower, S.~Rodriguez~Bernabeu, J.~Hammond, J.~Gunnels, ``Hunting for quantum advantage in electronic structure calculations is a highly non-trivial task,'' arXiv:2603.28648v1 (2026). \url{https://arxiv.org/abs/2603.28648}
\bibitem{rende2024} R.~Rende, L.~L.~Viteritti, L.~Bardone, F.~Becca, S.~Goldt, ``A simple linear algebra identity to optimize large-scale neural network quantum states,'' Commun.\ Phys.\ 7, 260 (2024); arXiv:2310.05715v2, Table~I. \url{https://arxiv.org/abs/2310.05715}
\bibitem{gu2026} Y.~Gu et al., ``Solving the Hubbard model with neural quantum states,'' arXiv:2507.02644v2 (2025). \url{https://arxiv.org/abs/2507.02644}
\bibitem{pfau2024} D.~Pfau, S.~Axelrod, H.~Sutterud, I.~von~Glehn, J.~S.~Spencer, ``Accurate computation of quantum excited states with neural networks,'' Science 385, eadn0137 (2024). \url{https://doi.org/10.1126/science.adn0137}
\bibitem{schatzle2026} Z.~Sch\"atzle et al., ``Ab initio simulation of excited-state potential energy surfaces with transferable deep quantum Monte Carlo,'' arXiv:2503.19847 (2025). \url{https://arxiv.org/abs/2503.19847}
\bibitem{layden2023} D.~Layden et al., ``Quantum-enhanced Markov chain Monte Carlo,'' Nature 619, 282--287 (2023). \url{https://doi.org/10.1038/s41586-023-06095-4}
\bibitem{orfi2024} A.~Orfi, D.~Sels, ``Bounding the speedup of the quantum-enhanced Markov-chain Monte Carlo algorithm,'' Phys.\ Rev.\ A 110, 052414 (2024). \url{https://doi.org/10.1103/PhysRevA.110.052414}
\bibitem{marshall2026} K.~V.~Marshall et al., ``Quantum-enhanced Markov chain Monte Carlo for combinatorial optimization,'' arXiv:2602.06171 (2026). \url{https://arxiv.org/abs/2602.06171}
\bibitem{hsieh2026} W.~T.~Hsieh, A.~Orfi, D.~Sels, ``Adiabatic dressing of quantum enhanced Markov chains,'' arXiv:2603.28076 (2026). \url{https://arxiv.org/abs/2603.28076}
\bibitem{yamamoto2022} A.~Yamamoto, ``Quantum sampling for the Euclidean path integral of lattice gauge theory,'' Phys.\ Rev.\ D 105, 094501 (2022). \url{https://doi.org/10.1103/PhysRevD.105.094501}
\bibitem{larose2026} R.~LaRose, ``A brief history of quantum vs classical computational advantage,'' Quantum 10, 2198 (2026). \url{https://doi.org/10.22331/q-2026-09-01-2198}
\bibitem{aa2014} S.~Aaronson, A.~Ambainis, ``The need for structure in quantum speedups,'' Theory of Computing 10, 133--166 (2014); arXiv:0911.0996. \url{https://arxiv.org/abs/0911.0996}
\bibitem{turturean2026} D.~Turturean, ``A proof of the Aaronson--Ambainis conjecture,'' preprint with Lean~4 formalization, repository commit \texttt{2a6ceae} (6 November 2026). \url{https://github.com/physical-superintelligence/aaronson-ambainis}
\bibitem{psi2026} Physical Superintelligence, ``A proof of the Aaronson--Ambainis conjecture'' (6 November 2026). \url{https://psi.inc/research/aaronson-ambainis/}
\bibitem{blancexpo2026} G.~Blanc et al., ``Proof of the Aaronson--Ambainis conjecture: an exposition'' (November 2026). \url{https://guyblanc.github.io/AA_Exposition.pdf}
\bibitem{aaronsonblog2026} S.~Aaronson, ``Aaronson--Ambainis conjecture proved! (By AI, of course),'' Shtetl-Optimized (November 2026). \url{https://scottaaronson.blog/?p=10295}
\bibitem{aaronson2019} S.~Aaronson, ``The Aaronson--Ambainis conjecture (2008--2019),'' Shtetl-Optimized (17 November 2019), on a proof announced that year.
\bibitem{bdst2026} G.~Blanc, J.~Docter, C.~Strassle, L.-Y.~Tan, ``Quantum speedups require structure or depth,'' arXiv:2608.19158 (2026); FOCS 2026. \url{https://arxiv.org/abs/2608.19158}
\bibitem{gmnr2026} U.~Girish, K.~Mittal, B.~Nehoran, R.~Raz, ``Distributional variants of the Aaronson--Ambainis conjecture,'' arXiv:2609.35327 (2026). \url{https://arxiv.org/abs/2609.35327}
\bibitem{beals2001} R.~Beals, H.~Buhrman, R.~Cleve, M.~Mosca, R.~de~Wolf, ``Quantum lower bounds by polynomials,'' J.\ ACM 48, 778--797 (2001). \url{https://doi.org/10.1145/502090.502097}
\bibitem{abkrt2021} S.~Aaronson, S.~Ben-David, R.~Kothari, S.~Rao, A.~Tal, ``Degree vs.\ approximate degree and quantum implications of Huang's sensitivity theorem,'' Proc.\ STOC 2021, 1330--1342. \url{https://doi.org/10.1145/3406325.3451047}
\bibitem{hhl2009} A.~W.~Harrow, A.~Hassidim, S.~Lloyd, ``Quantum algorithm for linear systems of equations,'' Phys.\ Rev.\ Lett.\ 103, 150502 (2009). \url{https://doi.org/10.1103/PhysRevLett.103.150502}
\bibitem{costa2022} P.~C.~S.~Costa, D.~An, Y.~R.~Sanders, Y.~Su, R.~Babbush, D.~W.~Berry, ``Optimal scaling quantum linear-systems solver via discrete adiabatic theorem,'' PRX Quantum 3, 040303 (2022). \url{https://doi.org/10.1103/PRXQuantum.3.040303}
\bibitem{clader2013} B.~D.~Clader, B.~C.~Jacobs, C.~R.~Sprouse, ``Preconditioned quantum linear system algorithm,'' Phys.\ Rev.\ Lett.\ 110, 250504 (2013). \url{https://doi.org/10.1103/PhysRevLett.110.250504}
\bibitem{montanaro2016} A.~Montanaro, S.~Pallister, ``Quantum algorithms and the finite element method,'' Phys.\ Rev.\ A 93, 032324 (2016); arXiv:1512.05903. \url{https://arxiv.org/abs/1512.05903}
\bibitem{scherer2017} A.~Scherer, B.~Valiron, S.-C.~Mau, S.~Alexander, E.~van den Berg, T.~E.~Chapuran, ``Concrete resource analysis of the quantum linear-system algorithm used to compute the electromagnetic scattering cross section of a 2D target,'' Quantum Inf.\ Process.\ 16, 60 (2017); arXiv:1505.06552v2. \url{https://arxiv.org/abs/1505.06552}
\bibitem{linden2022} N.~Linden, A.~Montanaro, C.~Shao, ``Quantum vs.\ classical algorithms for solving the heat equation,'' Commun.\ Math.\ Phys.\ 395, 601--641 (2022). \url{https://doi.org/10.1007/s00220-022-04442-6}
\bibitem{liu2021c} J.-P.~Liu, H.~{\O}.~Kolden, H.~K.~Krovi, N.~F.~Loureiro, K.~Trivisa, A.~M.~Childs, ``Efficient quantum algorithm for dissipative nonlinear differential equations,'' Proc.\ Natl.\ Acad.\ Sci.\ 118, e2026805118 (2021). \url{https://arxiv.org/abs/2011.03185}
\bibitem{lewis2024} D.~Lewis, S.~Eidenbenz, B.~Nadiga, Y.~Suba\c{s}{\i}, ``Limitations for quantum algorithms to solve turbulent and chaotic systems,'' Quantum 8, 1509 (2024). \url{https://doi.org/10.22331/q-2024-10-24-1509}
\bibitem{an2025} D.~An, J.-P.~Liu, D.~Wang, Q.~Zhao, ``Quantum differential equation solvers: limitations and fast-forwarding,'' Commun.\ Math.\ Phys.\ 406, 189 (2025); arXiv:2211.05246. \url{https://arxiv.org/abs/2211.05246}
\bibitem{penuel2025} J.~Penuel, A.~Katabarwa, P.~D.~Johnson, P.~Kuklinski, B.~Rempfer, C.~Farquhar, Y.~Cao, M.~C.~Garrett, ``Detailed assessment of calculating drag force with quantum computers: explicit time-evolution precludes exponential advantage for nonlinear differential equations,'' arXiv:2406.06323v3 (8 December 2025; first posted June 2024 as ``Feasibility of accelerating incompressible computational fluid dynamics simulations with fault-tolerant quantum computers''). \url{https://arxiv.org/abs/2406.06323}
\bibitem{li2025fluid} X.~Li, X.~Yin, N.~Wiebe, J.~Chun, G.~K.~Schenter, M.~S.~Cheung, J.~M\"ulmenst\"adt, ``Potential quantum advantage for simulation of fluid dynamics,'' Phys.\ Rev.\ Research 7, 013036 (2025). \url{https://arxiv.org/abs/2303.16550}
\bibitem{lapworth2022} L.~Lapworth, ``A hybrid quantum-classical CFD methodology with benchmark HHL solutions,'' arXiv:2206.00419 (2022). \url{https://arxiv.org/abs/2206.00419}
\bibitem{lapworth2025} L.~Lapworth, C.~S\"underhauf, ``Preconditioned block encodings for quantum linear systems,'' Quantum Sci.\ Technol.\ 10 (2025); arXiv:2502.20908. \url{https://arxiv.org/abs/2502.20908}
\bibitem{goldfriend2026} T.~Goldfriend, L.~Lapworth, N.~Yoran, A.~Naveh, ``Approximate quantum linear solvers for hybrid CFD: end-to-end analysis with a Chebyshev-LCU approach,'' arXiv:2606.01067 (2026). \url{https://arxiv.org/abs/2606.01067}
\bibitem{tu2025} Y.~Tu, M.~Dubynskyi, M.~Mohammadisiahroudi, E.~Riashchentceva, J.~Cheng, D.~Ryashchentsev, T.~Terlaky, J.~Liu, ``Towards identifying possible fault-tolerant advantage of quantum linear system algorithms in terms of space, time and energy,'' arXiv:2502.11239 (2025). \url{https://arxiv.org/abs/2502.11239}
\bibitem{cecile2026} M.~G.~Cecile, N.-L.~van H\"ulst, T.~Hashizume et al., ``Quantum algorithms for computational fluid dynamics,'' arXiv:2609.39962 (2026). \url{https://arxiv.org/abs/2609.39962}
\bibitem{dalzell2026} A.~M.~Dalzell, J.~Li, Y.~Su, ``Faster quantum linear system solver beyond the condition number,'' arXiv:2607.07691 (2026). \url{https://arxiv.org/abs/2607.07691}
\bibitem{babbush2023osc} R.~Babbush, D.~W.~Berry, R.~Kothari, R.~D.~Somma, N.~Wiebe, ``Exponential quantum speedup in simulating coupled classical oscillators,'' Phys.\ Rev.\ X 13, 041041 (2023). \url{https://arxiv.org/abs/2303.13012}
\bibitem{brandt1977} A.~Brandt, ``Multi-level adaptive solutions to boundary-value problems,'' Math.\ Comp.\ 31, 333--390 (1977). \url{https://doi.org/10.1090/S0025-5718-1977-0431719-X}
\bibitem{kochkov2021} D.~Kochkov, J.~A.~Smith, A.~Alieva, Q.~Wang, M.~P.~Brenner, S.~Hoyer, ``Machine learning--accelerated computational fluid dynamics,'' Proc.\ Natl.\ Acad.\ Sci.\ 118, e2101784118 (2021). \url{https://doi.org/10.1073/pnas.2101784118}
\bibitem{lam2023} R.~Lam et al., ``Learning skillful medium-range global weather forecasting,'' Science 382, 1416--1421 (2023). \url{https://doi.org/10.1126/science.adi2336}
\end{thebibliography}

\end{document}
